% !TeX root = ../main.tex
\section{Source Density and Positional Memory}

% TODO: add references on grafting/polarization experiments, including Kinneret-related experiments
% TODO: add references on source gradients
% TODO: define prescribed rho profiles
% TODO: define dynamic rho equation
% TODO: add simulations for grafted tissue configurations

\begin{itemize}
    \item Motivation from grafting/polarization experiments:
    \begin{itemize}
        \item Tissue fragments taken from different body regions can behave differently after being grafted.
        \item The outcome depends on how fragments are combined.
        \item This suggests positional memory.
    \end{itemize}
    \item Problem for classical Gierer--Meinhardt models:
    \begin{itemize}
        \item Away from the head, activator concentration may be very small.
        \item Classical Gierer--Meinhardt variables may not encode whether a fragment comes from the upper or lower body column.
        \item The model lacks a mechanism to remember tissue origin.
    \end{itemize}
    \item Source density/source gradient as positional information:
    \begin{itemize}
        \item \(\rho\) can encode distance from the head or positional identity.
        \item It is higher near the head and lower away from the head.
        \item It may act as a localizer for future head formation.
    \end{itemize}
    \item First modelling option: prescribed source density \(\rho(x)\).
    \begin{itemize}
        \item \(\rho(x)\) is fixed in time.
        \item It represents externally imposed positional information.
        \item Simulate grafted profiles by combining different pieces of \(\rho(x)\).
        \item Wounds are introduced near grafting points.
        \item Peaks can emerge near maxima or discontinuities/transition regions of \(\rho\).
        \item Head positions need not coincide with Turing-selected points.
    \end{itemize}
    \item Second modelling option: dynamic source density \(\rho(t,x)\).
    \begin{itemize}
        \item Initial condition \(\rho(0,x)\) encodes the grafted tissue configuration.
        \item \(\rho\) evolves slowly.
        \item \(\rho\) must react slower than activator-inhibitor dynamics to preserve memory.
        \item \(\rho\) must diffuse:
        \begin{itemize}
            \item not too fast, otherwise it becomes nearly constant;
            \item not too slow, otherwise it remains too localized near peaks;
            \item an intermediate diffusion rate can encode approximate distance from heads.
        \end{itemize}
    \end{itemize}
    \item Expected numerical behavior:
    \begin{itemize}
        \item initially similar to prescribed \(\rho(x)\);
        \item peaks form near source-density maxima;
        \item peaks can remain there for long times;
        \item this creates metastable positional memory;
        \item over longer times, peaks may drift toward Turing-selected patterns.
    \end{itemize}
    \item Main conclusion:
    \begin{itemize}
        \item Source density helps encode positional memory.
        \item However, dynamic source density alone may generate metastable transients rather than genuine stable stationary heads at arbitrary positions.
    \end{itemize}
\end{itemize}
